\section{Artifact and reproducibility}
\label{sec:reproducibility}

The complete artifact is available at
\url{https://github.com/AnanasClassic/megaminx_34_56_phases_reduction}
\cite{kuznetsov2026artifact}.
It contains the paper sources, compressed SQLite certificate databases, both
certificate-generation checkpoints for each pair, training and cascade logs, pinned
problem manifests, independent verifier source, and scripts for reconstructing
all omitted intermediate tables.  The bibliography pins the exact artifact
commit used for this manuscript; the repository landing-page URL is provided
only for navigation.  No DOI or third-party archival deposit is claimed at the
time of writing.

The artifact supports three levels of reproduction:
\begin{enumerate}
  \item \textbf{Package audit.}  \texttt{make verify-package} streams both
    compressed databases, verifies compressed and uncompressed hashes, checks
    their schemas and complete distributions, and verifies all checkpoint
    hashes.  It does not require a neural framework or regenerated tables.
  \item \textbf{Proof reproduction.}  \texttt{make verify-pair34} and
    \texttt{make verify-pair56} fetch the digest-pinned source snapshot,
    recompute the subgroup orders from the committed action manifests, rebuild
    the table generator and replay implementations, regenerate the phase
    tables, maximal records, physical products, equivalence analysis, and exact
    rewrites, and replay every certificate.  The neural checkpoints are not
    loaded.
  \item \textbf{Discovery reproduction.}  The pair-specific training and
    cascade scripts retrain a model or rerun certificate discovery on a GPU.
    Because neural training and parallel floating-point execution need not be
    bit-reproducible, success is judged only by the same deterministic
    certificate checker.
\end{enumerate}

The proof environment requires Python~3.10 or newer, SymPy for the exact
Schreier--Sims audit, and the digest-pinned Go~1.22.3 container image.  The
published discovery run used an NVIDIA A100 80GB PCIe GPU with CUDA/bfloat16
inference.  Proof replay is CPU-only.  The historical problem-manifest builder
also accepts a p900 generator file with a pinned SHA-256 digest; that file is
not redistributed here.  It is construction provenance rather than a
proof-time dependency: the committed manifests are checked directly against
the physical move implementation and their group orders are recomputed.

The two compressed databases contain 536,369 and 399,167 direct records.  The
committed reference reports are historical Python-only outputs whose schemas
contain no Go transcript fields.  A full current reproduction emits a new
report containing the Go binary digest and the exact cross-language transcript
digest.  The repository also documents the 108-byte
physical state format, strict FTM word grammar, coordinate conventions,
artifact binary formats, and an explicit corruption-oriented threat model.
Code and computational data are distributed under the repository's stated
CC~BY-NC~4.0 terms, matching the license of the source dataset on which the
coordinate reimplementation is based.
